% ECONOMICS_PROBLEM: {"id": "C78-134", "title": "Does double primality reduce the maximum number of stable matchings?", "jel_code": "C78", "source_id": "OP-0218"}
% FIRST_POSTED: 2026-10-09
% ATTEMPT_STATUS: solved
% ENTRY_KIND: research_attempt
\documentclass[11pt,a4paper]{article}
\usepackage[T1]{fontenc}
\usepackage[utf8]{inputenc}
\usepackage{lmodern,amsmath,amssymb,amsthm,mathtools}
\usepackage[margin=25mm]{geometry}
\usepackage{hyperref}
\hypersetup{colorlinks=true,urlcolor=blue,linkcolor=blue}
\setlength{\parindent}{0pt}
\setlength{\parskip}{5pt}
\newtheorem{theorem}{Theorem}
\newtheorem{proposition}[theorem]{Proposition}
\newtheorem{lemma}[theorem]{Lemma}
\newtheorem{corollary}[theorem]{Corollary}
\newcommand{\R}{\mathbb R}
\newcommand{\E}{\mathbb E}
\newcommand{\Prb}{\mathbb P}
\newcommand{\ind}{\mathbf 1}
\newcommand{\TV}{\operatorname{TV}}
\newcommand{\KL}{\operatorname{KL}}
\newcommand{\Var}{\operatorname{Var}}
\newcommand{\OPT}{\operatorname{OPT}}
\newcommand{\diam}{\operatorname{diam}}
\newcommand{\supp}{\operatorname{supp}}
\DeclareMathOperator*{\argmax}{arg\,max}
\DeclareMathOperator*{\argmin}{arg\,min}
\setlength{\emergencystretch}{3em}
\begin{document}
\noindent\textbf{Problem ID:} \texttt{C78-134}\quad\textbf{Model:} GPT 6 Astra Pro\quad\textbf{Version:} 2\par
\section*{Does double primality reduce the maximum number of stable matchings?}

\textbf{Status.} Disproved as written at $n=2$; the higher-order variant remains open. The displayed universal assertion in the canonical statement is false under its labelled reachable-event definition. We prove $S_2^{\mathrm{DP}}=1<2=S_2$. This is a negative resolution of the literal $n\geq1$, $n\ne4$ assertion, not a resolution of a question whose quantifier has been changed to $n\geq3$.

\textbf{Pro revision.} Version 2 retains the complete order-two proof from the original attempt and incorporates the supplied independent audit's scope clarification. The canonical definition and problem ID are unchanged. Neither the literal counterexample nor its completion estimate should be transferred to the amended higher-order question.

\subsection*{The exact assertion being tested}
In a complete strict stable-marriage instance, $n$ proposers and $n$ receivers each rank every member of the opposite side. A perfect matching is stable when no unmatched pair prefers each other to their assigned partners. Let $S_n$ be the largest number of stable matchings among all such instances.

Input decomposition means partitioning both sides into equally sized corresponding blocks such that every agent ranks every partner in its block above all outsiders. Input-prime means that there is no nontrivial such partition. For proposer-side serial deferred acceptance, label a proposal event $(p,j)$ for proposer $p$'s $(j+1)$st proposal. Let $\mathcal F$ consist of the event sets realized by all reachable states under all legal schedules, including the initial and terminal states; $E=\bigcup_{F\in\mathcal F}F$. Behavioural decomposition is a nontrivial partition $E=\bigsqcup E_j$ for which
\[
 \mathcal F=\left\{\bigsqcup_jF_j:
                  F_j\in\{F\cap E_j:F\in\mathcal F\}\right\}.
\]
An instance is doubly prime if it is prime for both notions. Define $S_n^{\mathrm{DP}}$ by maximizing stable-set size over doubly prime instances. The catalogue explicitly asks whether
\[
 \forall n\geq1\quad(n\ne4\ \Longrightarrow\ S_n^{\mathrm{DP}}=S_n).\tag{1}
\]
In particular, (1) includes $n=2$, and all following arguments use exactly this proposing side and event-family convention.

\begin{lemma}[Two stable matchings force distinct first choices at order two]
If a complete strict $2\times2$ instance has two stable matchings, the proposers' first choices are distinct. Its reachable event family is the full power set of its two first-proposal events, so it is behaviourally decomposable.
\end{lemma}
\begin{proof}
Write the proposers as $p,q$ and receivers as $a,b$. There are exactly two possible perfect matchings,
$M=\{(p,a),(q,b)\}$ and $N=\{(p,b),(q,a)\}$.
Suppose both proposers prefer $a$ to $b$. Stability of $M$ then requires $a$ to prefer $p$ to $q$, otherwise $(q,a)$ blocks. Stability of $N$ requires $a$ to prefer $q$ to $p$, otherwise $(p,a)$ blocks. Strict preferences cannot satisfy both. The same argument applies if both proposers prefer $b$. Thus their first choices are distinct.

With distinct first choices, neither proposer ever faces competition for its first receiver. Both first proposals are accepted and no rejection occurs. Each proposer makes exactly one proposal. Either proposal can be first in a legal schedule, and the process terminates after both. Hence, writing the events as $A,B$,
$\mathcal F=\{\varnothing,\{A\},\{B\},\{A,B\}\}$. This is the product of its projections on $\{A\}$ and $\{B\}$, a nontrivial behavioural decomposition.
\end{proof}

\begin{proposition}[The unrestricted maximum at order two]
$S_2=2$.
\end{proposition}
\begin{proof}
There are only two perfect matchings, so $S_2\leq2$. In the instance
\[
 p:a\succ b,\quad q:b\succ a,\qquad
 a:q\succ p,\quad b:p\succ q,
\]
$M=\{(p,a),(q,b)\}$ is stable because both proposers have their first choices, while $N=\{(p,b),(q,a)\}$ is stable because both receivers have their first choices. Thus two stable matchings are attainable.
\end{proof}

\begin{proposition}[A doubly prime witness with one stable matching]
Consider
\[
 p:a\succ b,\qquad q:a\succ b,\qquad
 a:p\succ q,\qquad b:q\succ p.                         \tag{2}
\]
This instance is input-prime and behaviourally prime, and has exactly one stable matching.
\end{proposition}
\begin{proof}
At order two a nontrivial input partition must consist of two singleton pairs, each mutually ranked first. Instance (2) has the mutual first pair $(p,a)$, but its complement $(q,b)$ is not mutual first, since $q$ prefers $a$. The alternative singleton pairing also fails. Thus no nontrivial input partition exists.

The unique stable matching is $M=\{(p,a),(q,b)\}$: it is stable directly, and the alternative matching is blocked by the mutual first pair $(p,a)$.

For the reachable events, put $A=(p,0)$, $B=(q,0)$ and $C=(q,1)$. Proposer $p$ makes exactly one proposal, always held by $a$. Proposer $q$ first proposes to $a$, and can next propose to $b$ exactly after both first proposals have occurred. The complete reachable family is
\[
 \mathcal F=\{\varnothing,\{A\},\{B\},\{A,B\},\{A,B,C\}\}.\tag{3}
\]
Both orders of the first two proposals are legal, verifying attainability of every listed set. Event $C$ needs both $A$ and $B$, verifying that no other set is reachable.
Suppose (3) factored over a nontrivial partition of $\{A,B,C\}$. The factor containing $C$ must also contain $A$: otherwise combine that factor's projection of the full set with the empty projection on the factor containing $A$, and any permitted projections on remaining factors. The resulting union contains $C$ without $A$, contradicting (3). The identical argument forces $B$ into the factor containing $C$. All three events would be in one factor, contradicting a nontrivial partition. Hence the event family is prime.
\end{proof}

\begin{theorem}[Failure of the literal universal claim]
Under the canonical definitions, $S_2^{\mathrm{DP}}=1$ and $S_2=2$. Therefore assertion (1) is false.
\end{theorem}
\begin{proof}
The lemma rules out two stable matchings in every behaviourally prime order-two instance, giving $S_2^{\mathrm{DP}}\leq1$. Proposition 3 gives a doubly prime instance attaining one. Proposition 2 gives the unrestricted maximum. Since $2\ne4$, the inequality at $n=2$ contradicts (1).
\end{proof}

\subsection*{Finite verification and source qualification}
The supplied independent audit, \emph{Independent checks for the economics ``solved'' audit}, uses \texttt{verify\_matching.py} to enumerate all $(2!)^4=16$ labelled strict order-two preference profiles. It computes their stable matchings, explores every reachable serial deferred-acceptance state, and tests every nontrivial proposer bipartition for a product of the counter families. The supplied \texttt{matching\_results.json} reports $S_2=2$, $S_2^{\mathrm{DP}}=1$, and eight behaviourally prime profiles. These values agree with the proofs above.

The reduction from event partitions to proposer partitions uses the prefix property, not a graph heuristic. If two events of one proposer were in different factors, the projection of a reachable set containing the later event could be combined with the empty projection of the earlier event's factor, producing a later event without its required predecessor. Thus a factor cannot split a proposer's events. At order two, distinct first choices give the product family in Lemma 1; every behaviourally prime profile therefore has coincident first choices. Such a profile cannot admit a two-singleton input-independent partition. Hence the audit's behaviourally prime class is indeed the doubly prime class here.

The exhaustive finite check establishes only the order-two values. It does not enumerate larger orders, prove the higher-order extremal conjecture, or constitute proof-assistant formalization. The mathematical counterexample above does not depend on that computation.

The primary source's Section 7, question E, discusses stable-set maxima starting at $n=3$, then $4,5,6$, and asks whether order four is the only cost. Its preceding discussion does include doubly prime instances of order two, but question E does not give an order-two stable-set count. The present counterexample concerns the explicit $n\geq1$ quantifier in the canonical formulation. An amended higher-order question restricted to $n\geq3$, $n\ne4$ is not answered here; neither are exact maxima at larger unproved orders. This small-order issue should not be presented as a proof of the large-order extremal conjecture or a general method of constructing extremizers.

\section{Completion Estimate}
100\%

This is a post hoc, heuristic estimate for the literal full catalogue target.
The counterexample proves the canonical universal assertion false by establishing the exact strict inequality at order two. This estimate records that negative resolution only; it makes no percentage claim for the amended $n\geq3$, $n\ne4$ question, which remains open here.

\subsection*{Reference}
Yoshiteru Ishida (2026), \emph{Prime factorisation of stable-matching instances: uniqueness, simultaneous products, and an exact census}, arXiv:2610.05617v1, Sections 2 and 4 and Section 7, question E. \url{https://arxiv.org/html/2610.05617v1}. The labelled-event definition and the text of question E were checked in the primary version. The order-two argument is supplied in full above.

\emph{Independent checks for the economics ``solved'' audit} (9 October 2026), user-supplied audit attachments \texttt{README.md}, \texttt{proposed\_corrections.md}, \texttt{verify\_matching.py}, and \texttt{matching\_results.json}.

The reviewed repository snapshot is \texttt{ccccd798864807c1f3fb8fe36727fe35ee875aeb}. The script uses Python's standard library, and its finite output is supplementary evidence rather than an external universal theorem.

\end{document}
